% =====================================================================
%  CV Etienne POUILLE — une colonne, texte sélectionnable, compatible ATS
%  Compilation (XeLaTeX) :
%    Standard : xelatex -jobname=CV_Etienne_POUILLE cv.tex
%    Ciblé    : xelatex -jobname=CV_Etienne_POUILLE \
%               '\def\cible{Stage AI Engineer}% =====================================================================
%  CV Etienne POUILLE — une colonne, texte sélectionnable, compatible ATS
%  Compilation (XeLaTeX) :
%    Standard : xelatex -jobname=CV_Etienne_POUILLE cv.tex
%    Ciblé    : xelatex -jobname=CV_Etienne_POUILLE \
%               '\def\cible{Stage AI Engineer}% =====================================================================
%  CV Etienne POUILLE — une colonne, texte sélectionnable, compatible ATS
%  Compilation (XeLaTeX) :
%    Standard : xelatex -jobname=CV_Etienne_POUILLE cv.tex
%    Ciblé    : xelatex -jobname=CV_Etienne_POUILLE \
%               '\def\cible{Stage AI Engineer}% =====================================================================
%  CV Etienne POUILLE — une colonne, texte sélectionnable, compatible ATS
%  Compilation (XeLaTeX) :
%    Standard : xelatex -jobname=CV_Etienne_POUILLE cv.tex
%    Ciblé    : xelatex -jobname=CV_Etienne_POUILLE \
%               '\def\cible{Stage AI Engineer}\input{cv.tex}'
%  Sur Overleaf (compilateur : XeLaTeX) : décommenter la ligne ci-dessous
%  pour obtenir la version ciblée, et y modifier l'intitulé du poste.
% =====================================================================
%\def\cible{Stage AI Engineer}
\documentclass[9pt,a4paper]{extarticle}
\usepackage[top=1.5cm,bottom=1.3cm,left=2cm,right=2cm]{geometry}
\usepackage{fontspec}
\usepackage{xcolor}
\usepackage{enumitem}
\usepackage[hidelinks,pdfusetitle]{hyperref}

% --- Polices : Inter si présente (graisse SemiBold pour le gras), sinon TeX Gyre Heros
\IfFontExistsTF{Inter SemiBold}
  {\setmainfont{Inter}[Ligatures=NoCommon,BoldFont={Inter SemiBold},RawFeature={-calt;-case}]}
  {\setmainfont{TeX Gyre Heros}[Ligatures=NoCommon]}

% --- Couleurs (vert foncé du portfolio)
\definecolor{ink}{HTML}{1B2922}
\definecolor{accent}{HTML}{2F5A44}
\definecolor{muted}{HTML}{55685D}
\definecolor{rule}{HTML}{B8C6BD}

\hypersetup{pdftitle={CV - Etienne Pouille},pdfauthor={Etienne Pouille},
            pdfsubject={CV - Stage IA / MLOps / Edge AI / DevOps}}

\pagestyle{empty}
\setlength{\parindent}{0pt}
\setlength{\parskip}{0pt}
\raggedright
\AtBeginDocument{\color{ink}}
\setlist[itemize]{leftmargin=1em,itemsep=1.5pt,topsep=2pt,parsep=0pt,label={•}}
\linespread{1.16}

% --- Macros
% Titre de section : petit, majuscules, filet discret
\newcommand{\sect}[1]{\par\vspace{13pt}%
  {\bfseries\color{accent}\small\addfontfeature{LetterSpace=5}\MakeUppercase{#1}}\par\vspace{-2pt}%
  {\color{rule}\rule{\linewidth}{0.4pt}}\par\vspace{4pt}}
% Ligne « titre – précision ........ date »
\newcommand{\entry}[3]{\par\noindent{\bfseries #1}{\color{muted}\;·\; #2}\hfill{\color{muted}#3}\par\nopagebreak}
% Projet : nom en gras puis une phrase
\newcommand{\projet}[2]{\par\vspace{4.5pt}\noindent{\bfseries #1}{\color{muted}\;–}\;#2\par}
\newcommand{\skill}[2]{\par\noindent{\bfseries #1}~:\;{\color{muted}#2}\par}

\begin{document}

% --- En-tête (centré)
\begingroup\centering
\ifdefined\cible
  {\bfseries\color{accent}\fontsize{21}{24}\selectfont \cible\par}
  \vspace{3pt}
  {\bfseries\large Etienne POUILLE}\par
  \vspace{1pt}
  {\color{muted}Étudiant en 3e année à EPITECH Lille · stage du 29 mars au 27 août 2027 · Nord de la France}\par
\else
  {\bfseries\color{accent}\fontsize{24}{27}\selectfont Etienne POUILLE\par}
  \vspace{3pt}
  {\large Étudiant en 3e année à EPITECH Lille}\par
  \vspace{1pt}
  {\color{muted}AI Engineer · MLOps · Edge AI · DevOps\; —\; stage du 29 mars au 27 août 2027 · Nord de la France}\par
\fi
\vspace{5pt}
{\small
\href{https://etienne-pouille.work}{\bfseries\color{accent}Portfolio : etienne-pouille.work}{\color{muted}\; · \;\href{mailto:contact@etienne-pouille.work}{contact@etienne-pouille.work} · 06~24~37~79~40}\par}
\endgroup

\vspace{9pt}
Je construis des systèmes d’IA et l’infrastructure qui les fait tourner : agents et LLM locaux (Ollama),
IA embarquée sur voiture autonome, pipelines CI/CD.

% --- Expérience
\sect{Expérience}
\entry{Stage Ingénieur SysOps / DevOps}{Cyllene}{sept. – déc. 2025}
\begin{itemize}
  \item Outil Python de génération et de visualisation automatiques des infrastructures clients,
        développé en autonomie et intégré aux pipelines GitLab CI/CD.
  \item Utilisé pour faciliter les échanges techniques avec les clients et les équipes comptabilité et développement.
\end{itemize}

% --- Projets
\sect{Projets}
\projet{Noctia}{Plugin Claude Code : règles de développement et modes de travail (l’IA en mentor, ou en autonomie avec orchestrateur et sous-agents).}
\projet{ALF}{Moteur agentique en Python connecté aux APIs Ollama, avec interface en ligne de commande.}
\projet{Tri d’e-mails}{Tri automatique par LLM local (API Ollama) avec n8n, déployé sur VPS personnel.}
\projet{RoboCar}{Voiture quasi-autonome (EPITECH, équipe de 5) : suivi de ligne blanche, puis LiDAR et caméra 3D pour une conduite autonome en tout environnement.}
\projet{MyTorch}{Réseau de neurones recodé from scratch, appliqué au jeu d’échecs (EPITECH, en cours).}
\projet{StartTrek}{Reinforcement Learning pour l’alunissage autonome d’un module spatial simulé (EPITECH).}

% --- Compétences
\sect{Compétences}
\skill{IA et MLOps}{Python, Ollama (LLM locaux), agents, Claude Code, Reinforcement Learning}
\vspace{2pt}
\skill{DevOps}{Linux, Docker, GitLab CI/CD, GitHub Actions, Jenkins, Tailscale, n8n, Kubernetes (en apprentissage)}
\vspace{2pt}
\skill{Systèmes}{C, Rust, C++, Haskell, Bash}
\vspace{2pt}
\skill{Web}{React, JavaScript, HTML/CSS}

% --- Formation
\sect{Formation}
\entry{EPITECH Lille}{Programme Grande École en informatique, 5 ans}{2024 – présent}
{\color{muted}1re et 2e années validées avec les félicitations · spécialisation visée : IA, Maths, DevOps\\
Badge Google DeepMind « AI Research Foundations » (en cours)\par}
\vspace{4pt}
\entry{Baccalauréat Général, mention Bien}{Lycée Notre-Dame des Dunes, Dunkerque}{2021 – 2024}
{\color{muted}Section européenne anglais · spécialités NSI, Mathématiques, Maths expertes\par}

% --- Langues et intérêts
\sect{Langues et centres d’intérêt}
Français (natif) · Anglais B2–C1 (TEPitech/TOEIC 910, Cambridge FIRST) · Espagnol A2\par
\vspace{2pt}
{\color{muted}Escalade, callisthénie, sport automobile, musique, dessin\par}
\vspace{2pt}
{\color{muted}Séjours : Chicago et Mexico (familles d’accueil, 2022-2023), New York (2024), Londres, Paris, Bruxelles\par}

\end{document}
'
%  Sur Overleaf (compilateur : XeLaTeX) : décommenter la ligne ci-dessous
%  pour obtenir la version ciblée, et y modifier l'intitulé du poste.
% =====================================================================
%\def\cible{Stage AI Engineer}
\documentclass[9pt,a4paper]{extarticle}
\usepackage[top=1.5cm,bottom=1.3cm,left=2cm,right=2cm]{geometry}
\usepackage{fontspec}
\usepackage{xcolor}
\usepackage{enumitem}
\usepackage[hidelinks,pdfusetitle]{hyperref}

% --- Polices : Inter si présente (graisse SemiBold pour le gras), sinon TeX Gyre Heros
\IfFontExistsTF{Inter SemiBold}
  {\setmainfont{Inter}[Ligatures=NoCommon,BoldFont={Inter SemiBold},RawFeature={-calt;-case}]}
  {\setmainfont{TeX Gyre Heros}[Ligatures=NoCommon]}

% --- Couleurs (vert foncé du portfolio)
\definecolor{ink}{HTML}{1B2922}
\definecolor{accent}{HTML}{2F5A44}
\definecolor{muted}{HTML}{55685D}
\definecolor{rule}{HTML}{B8C6BD}

\hypersetup{pdftitle={CV - Etienne Pouille},pdfauthor={Etienne Pouille},
            pdfsubject={CV - Stage IA / MLOps / Edge AI / DevOps}}

\pagestyle{empty}
\setlength{\parindent}{0pt}
\setlength{\parskip}{0pt}
\raggedright
\AtBeginDocument{\color{ink}}
\setlist[itemize]{leftmargin=1em,itemsep=1.5pt,topsep=2pt,parsep=0pt,label={•}}
\linespread{1.16}

% --- Macros
% Titre de section : petit, majuscules, filet discret
\newcommand{\sect}[1]{\par\vspace{13pt}%
  {\bfseries\color{accent}\small\addfontfeature{LetterSpace=5}\MakeUppercase{#1}}\par\vspace{-2pt}%
  {\color{rule}\rule{\linewidth}{0.4pt}}\par\vspace{4pt}}
% Ligne « titre – précision ........ date »
\newcommand{\entry}[3]{\par\noindent{\bfseries #1}{\color{muted}\;·\; #2}\hfill{\color{muted}#3}\par\nopagebreak}
% Projet : nom en gras puis une phrase
\newcommand{\projet}[2]{\par\vspace{4.5pt}\noindent{\bfseries #1}{\color{muted}\;–}\;#2\par}
\newcommand{\skill}[2]{\par\noindent{\bfseries #1}~:\;{\color{muted}#2}\par}

\begin{document}

% --- En-tête (centré)
\begingroup\centering
\ifdefined\cible
  {\bfseries\color{accent}\fontsize{21}{24}\selectfont \cible\par}
  \vspace{3pt}
  {\bfseries\large Etienne POUILLE}\par
  \vspace{1pt}
  {\color{muted}Étudiant en 3e année à EPITECH Lille · stage du 29 mars au 27 août 2027 · Nord de la France}\par
\else
  {\bfseries\color{accent}\fontsize{24}{27}\selectfont Etienne POUILLE\par}
  \vspace{3pt}
  {\large Étudiant en 3e année à EPITECH Lille}\par
  \vspace{1pt}
  {\color{muted}AI Engineer · MLOps · Edge AI · DevOps\; —\; stage du 29 mars au 27 août 2027 · Nord de la France}\par
\fi
\vspace{5pt}
{\small
\href{https://etienne-pouille.work}{\bfseries\color{accent}Portfolio : etienne-pouille.work}{\color{muted}\; · \;\href{mailto:contact@etienne-pouille.work}{contact@etienne-pouille.work} · 06~24~37~79~40}\par}
\endgroup

\vspace{9pt}
Je construis des systèmes d’IA et l’infrastructure qui les fait tourner : agents et LLM locaux (Ollama),
IA embarquée sur voiture autonome, pipelines CI/CD.

% --- Expérience
\sect{Expérience}
\entry{Stage Ingénieur SysOps / DevOps}{Cyllene}{sept. – déc. 2025}
\begin{itemize}
  \item Outil Python de génération et de visualisation automatiques des infrastructures clients,
        développé en autonomie et intégré aux pipelines GitLab CI/CD.
  \item Utilisé pour faciliter les échanges techniques avec les clients et les équipes comptabilité et développement.
\end{itemize}

% --- Projets
\sect{Projets}
\projet{Noctia}{Plugin Claude Code : règles de développement et modes de travail (l’IA en mentor, ou en autonomie avec orchestrateur et sous-agents).}
\projet{ALF}{Moteur agentique en Python connecté aux APIs Ollama, avec interface en ligne de commande.}
\projet{Tri d’e-mails}{Tri automatique par LLM local (API Ollama) avec n8n, déployé sur VPS personnel.}
\projet{RoboCar}{Voiture quasi-autonome (EPITECH, équipe de 5) : suivi de ligne blanche, puis LiDAR et caméra 3D pour une conduite autonome en tout environnement.}
\projet{MyTorch}{Réseau de neurones recodé from scratch, appliqué au jeu d’échecs (EPITECH, en cours).}
\projet{StartTrek}{Reinforcement Learning pour l’alunissage autonome d’un module spatial simulé (EPITECH).}

% --- Compétences
\sect{Compétences}
\skill{IA et MLOps}{Python, Ollama (LLM locaux), agents, Claude Code, Reinforcement Learning}
\vspace{2pt}
\skill{DevOps}{Linux, Docker, GitLab CI/CD, GitHub Actions, Jenkins, Tailscale, n8n, Kubernetes (en apprentissage)}
\vspace{2pt}
\skill{Systèmes}{C, Rust, C++, Haskell, Bash}
\vspace{2pt}
\skill{Web}{React, JavaScript, HTML/CSS}

% --- Formation
\sect{Formation}
\entry{EPITECH Lille}{Programme Grande École en informatique, 5 ans}{2024 – présent}
{\color{muted}1re et 2e années validées avec les félicitations · spécialisation visée : IA, Maths, DevOps\\
Badge Google DeepMind « AI Research Foundations » (en cours)\par}
\vspace{4pt}
\entry{Baccalauréat Général, mention Bien}{Lycée Notre-Dame des Dunes, Dunkerque}{2021 – 2024}
{\color{muted}Section européenne anglais · spécialités NSI, Mathématiques, Maths expertes\par}

% --- Langues et intérêts
\sect{Langues et centres d’intérêt}
Français (natif) · Anglais B2–C1 (TEPitech/TOEIC 910, Cambridge FIRST) · Espagnol A2\par
\vspace{2pt}
{\color{muted}Escalade, callisthénie, sport automobile, musique, dessin\par}
\vspace{2pt}
{\color{muted}Séjours : Chicago et Mexico (familles d’accueil, 2022-2023), New York (2024), Londres, Paris, Bruxelles\par}

\end{document}
'
%  Sur Overleaf (compilateur : XeLaTeX) : décommenter la ligne ci-dessous
%  pour obtenir la version ciblée, et y modifier l'intitulé du poste.
% =====================================================================
%\def\cible{Stage AI Engineer}
\documentclass[9pt,a4paper]{extarticle}
\usepackage[top=1.5cm,bottom=1.3cm,left=2cm,right=2cm]{geometry}
\usepackage{fontspec}
\usepackage{xcolor}
\usepackage{enumitem}
\usepackage[hidelinks,pdfusetitle]{hyperref}

% --- Polices : Inter si présente (graisse SemiBold pour le gras), sinon TeX Gyre Heros
\IfFontExistsTF{Inter SemiBold}
  {\setmainfont{Inter}[Ligatures=NoCommon,BoldFont={Inter SemiBold},RawFeature={-calt;-case}]}
  {\setmainfont{TeX Gyre Heros}[Ligatures=NoCommon]}

% --- Couleurs (vert foncé du portfolio)
\definecolor{ink}{HTML}{1B2922}
\definecolor{accent}{HTML}{2F5A44}
\definecolor{muted}{HTML}{55685D}
\definecolor{rule}{HTML}{B8C6BD}

\hypersetup{pdftitle={CV - Etienne Pouille},pdfauthor={Etienne Pouille},
            pdfsubject={CV - Stage IA / MLOps / Edge AI / DevOps}}

\pagestyle{empty}
\setlength{\parindent}{0pt}
\setlength{\parskip}{0pt}
\raggedright
\AtBeginDocument{\color{ink}}
\setlist[itemize]{leftmargin=1em,itemsep=1.5pt,topsep=2pt,parsep=0pt,label={•}}
\linespread{1.16}

% --- Macros
% Titre de section : petit, majuscules, filet discret
\newcommand{\sect}[1]{\par\vspace{13pt}%
  {\bfseries\color{accent}\small\addfontfeature{LetterSpace=5}\MakeUppercase{#1}}\par\vspace{-2pt}%
  {\color{rule}\rule{\linewidth}{0.4pt}}\par\vspace{4pt}}
% Ligne « titre – précision ........ date »
\newcommand{\entry}[3]{\par\noindent{\bfseries #1}{\color{muted}\;·\; #2}\hfill{\color{muted}#3}\par\nopagebreak}
% Projet : nom en gras puis une phrase
\newcommand{\projet}[2]{\par\vspace{4.5pt}\noindent{\bfseries #1}{\color{muted}\;–}\;#2\par}
\newcommand{\skill}[2]{\par\noindent{\bfseries #1}~:\;{\color{muted}#2}\par}

\begin{document}

% --- En-tête (centré)
\begingroup\centering
\ifdefined\cible
  {\bfseries\color{accent}\fontsize{21}{24}\selectfont \cible\par}
  \vspace{3pt}
  {\bfseries\large Etienne POUILLE}\par
  \vspace{1pt}
  {\color{muted}Étudiant en 3e année à EPITECH Lille · stage du 29 mars au 27 août 2027 · Nord de la France}\par
\else
  {\bfseries\color{accent}\fontsize{24}{27}\selectfont Etienne POUILLE\par}
  \vspace{3pt}
  {\large Étudiant en 3e année à EPITECH Lille}\par
  \vspace{1pt}
  {\color{muted}AI Engineer · MLOps · Edge AI · DevOps\; —\; stage du 29 mars au 27 août 2027 · Nord de la France}\par
\fi
\vspace{5pt}
{\small
\href{https://etienne-pouille.work}{\bfseries\color{accent}Portfolio : etienne-pouille.work}{\color{muted}\; · \;\href{mailto:contact@etienne-pouille.work}{contact@etienne-pouille.work} · 06~24~37~79~40}\par}
\endgroup

\vspace{9pt}
Je construis des systèmes d’IA et l’infrastructure qui les fait tourner : agents et LLM locaux (Ollama),
IA embarquée sur voiture autonome, pipelines CI/CD.

% --- Expérience
\sect{Expérience}
\entry{Stage Ingénieur SysOps / DevOps}{Cyllene}{sept. – déc. 2025}
\begin{itemize}
  \item Outil Python de génération et de visualisation automatiques des infrastructures clients,
        développé en autonomie et intégré aux pipelines GitLab CI/CD.
  \item Utilisé pour faciliter les échanges techniques avec les clients et les équipes comptabilité et développement.
\end{itemize}

% --- Projets
\sect{Projets}
\projet{Noctia}{Plugin Claude Code : règles de développement et modes de travail (l’IA en mentor, ou en autonomie avec orchestrateur et sous-agents).}
\projet{ALF}{Moteur agentique en Python connecté aux APIs Ollama, avec interface en ligne de commande.}
\projet{Tri d’e-mails}{Tri automatique par LLM local (API Ollama) avec n8n, déployé sur VPS personnel.}
\projet{RoboCar}{Voiture quasi-autonome (EPITECH, équipe de 5) : suivi de ligne blanche, puis LiDAR et caméra 3D pour une conduite autonome en tout environnement.}
\projet{MyTorch}{Réseau de neurones recodé from scratch, appliqué au jeu d’échecs (EPITECH, en cours).}
\projet{StartTrek}{Reinforcement Learning pour l’alunissage autonome d’un module spatial simulé (EPITECH).}

% --- Compétences
\sect{Compétences}
\skill{IA et MLOps}{Python, Ollama (LLM locaux), agents, Claude Code, Reinforcement Learning}
\vspace{2pt}
\skill{DevOps}{Linux, Docker, GitLab CI/CD, GitHub Actions, Jenkins, Tailscale, n8n, Kubernetes (en apprentissage)}
\vspace{2pt}
\skill{Systèmes}{C, Rust, C++, Haskell, Bash}
\vspace{2pt}
\skill{Web}{React, JavaScript, HTML/CSS}

% --- Formation
\sect{Formation}
\entry{EPITECH Lille}{Programme Grande École en informatique, 5 ans}{2024 – présent}
{\color{muted}1re et 2e années validées avec les félicitations · spécialisation visée : IA, Maths, DevOps\\
Badge Google DeepMind « AI Research Foundations » (en cours)\par}
\vspace{4pt}
\entry{Baccalauréat Général, mention Bien}{Lycée Notre-Dame des Dunes, Dunkerque}{2021 – 2024}
{\color{muted}Section européenne anglais · spécialités NSI, Mathématiques, Maths expertes\par}

% --- Langues et intérêts
\sect{Langues et centres d’intérêt}
Français (natif) · Anglais B2–C1 (TEPitech/TOEIC 910, Cambridge FIRST) · Espagnol A2\par
\vspace{2pt}
{\color{muted}Escalade, callisthénie, sport automobile, musique, dessin\par}
\vspace{2pt}
{\color{muted}Séjours : Chicago et Mexico (familles d’accueil, 2022-2023), New York (2024), Londres, Paris, Bruxelles\par}

\end{document}
'
%  Sur Overleaf (compilateur : XeLaTeX) : décommenter la ligne ci-dessous
%  pour obtenir la version ciblée, et y modifier l'intitulé du poste.
% =====================================================================
%\def\cible{Stage AI Engineer}
\documentclass[9pt,a4paper]{extarticle}
\usepackage[top=1.5cm,bottom=1.3cm,left=2cm,right=2cm]{geometry}
\usepackage{fontspec}
\usepackage{xcolor}
\usepackage{enumitem}
\usepackage[hidelinks,pdfusetitle]{hyperref}

% --- Polices : Inter si présente (graisse SemiBold pour le gras), sinon TeX Gyre Heros
\IfFontExistsTF{Inter SemiBold}
  {\setmainfont{Inter}[Ligatures=NoCommon,BoldFont={Inter SemiBold},RawFeature={-calt;-case}]}
  {\setmainfont{TeX Gyre Heros}[Ligatures=NoCommon]}

% --- Couleurs (vert foncé du portfolio)
\definecolor{ink}{HTML}{1B2922}
\definecolor{accent}{HTML}{2F5A44}
\definecolor{muted}{HTML}{55685D}
\definecolor{rule}{HTML}{B8C6BD}

\hypersetup{pdftitle={CV - Etienne Pouille},pdfauthor={Etienne Pouille},
            pdfsubject={CV - Stage IA / MLOps / Edge AI / DevOps}}

\pagestyle{empty}
\setlength{\parindent}{0pt}
\setlength{\parskip}{0pt}
\raggedright
\AtBeginDocument{\color{ink}}
\setlist[itemize]{leftmargin=1em,itemsep=1.5pt,topsep=2pt,parsep=0pt,label={•}}
\linespread{1.16}

% --- Macros
% Titre de section : petit, majuscules, filet discret
\newcommand{\sect}[1]{\par\vspace{13pt}%
  {\bfseries\color{accent}\small\addfontfeature{LetterSpace=5}\MakeUppercase{#1}}\par\vspace{-2pt}%
  {\color{rule}\rule{\linewidth}{0.4pt}}\par\vspace{4pt}}
% Ligne « titre – précision ........ date »
\newcommand{\entry}[3]{\par\noindent{\bfseries #1}{\color{muted}\;·\; #2}\hfill{\color{muted}#3}\par\nopagebreak}
% Projet : nom en gras puis une phrase
\newcommand{\projet}[2]{\par\vspace{4.5pt}\noindent{\bfseries #1}{\color{muted}\;–}\;#2\par}
\newcommand{\skill}[2]{\par\noindent{\bfseries #1}~:\;{\color{muted}#2}\par}

\begin{document}

% --- En-tête (centré)
\begingroup\centering
\ifdefined\cible
  {\bfseries\color{accent}\fontsize{21}{24}\selectfont \cible\par}
  \vspace{3pt}
  {\bfseries\large Etienne POUILLE}\par
  \vspace{1pt}
  {\color{muted}Étudiant en 3e année à EPITECH Lille · stage du 29 mars au 27 août 2027 · Nord de la France}\par
\else
  {\bfseries\color{accent}\fontsize{24}{27}\selectfont Etienne POUILLE\par}
  \vspace{3pt}
  {\large Étudiant en 3e année à EPITECH Lille}\par
  \vspace{1pt}
  {\color{muted}AI Engineer · MLOps · Edge AI · DevOps\; —\; stage du 29 mars au 27 août 2027 · Nord de la France}\par
\fi
\vspace{5pt}
{\small
\href{https://etienne-pouille.work}{\bfseries\color{accent}Portfolio : etienne-pouille.work}{\color{muted}\; · \;\href{mailto:contact@etienne-pouille.work}{contact@etienne-pouille.work} · 06~24~37~79~40}\par}
\endgroup

\vspace{9pt}
Je construis des systèmes d’IA et l’infrastructure qui les fait tourner : agents et LLM locaux (Ollama),
IA embarquée sur voiture autonome, pipelines CI/CD.

% --- Expérience
\sect{Expérience}
\entry{Stage Ingénieur SysOps / DevOps}{Cyllene}{sept. – déc. 2025}
\begin{itemize}
  \item Outil Python de génération et de visualisation automatiques des infrastructures clients,
        développé en autonomie et intégré aux pipelines GitLab CI/CD.
  \item Utilisé pour faciliter les échanges techniques avec les clients et les équipes comptabilité et développement.
\end{itemize}

% --- Projets
\sect{Projets}
\projet{Noctia}{Plugin Claude Code : règles de développement et modes de travail (l’IA en mentor, ou en autonomie avec orchestrateur et sous-agents).}
\projet{ALF}{Moteur agentique en Python connecté aux APIs Ollama, avec interface en ligne de commande.}
\projet{Tri d’e-mails}{Tri automatique par LLM local (API Ollama) avec n8n, déployé sur VPS personnel.}
\projet{RoboCar}{Voiture quasi-autonome (EPITECH, équipe de 5) : suivi de ligne blanche, puis LiDAR et caméra 3D pour une conduite autonome en tout environnement.}
\projet{MyTorch}{Réseau de neurones recodé from scratch, appliqué au jeu d’échecs (EPITECH, en cours).}
\projet{StartTrek}{Reinforcement Learning pour l’alunissage autonome d’un module spatial simulé (EPITECH).}

% --- Compétences
\sect{Compétences}
\skill{IA et MLOps}{Python, Ollama (LLM locaux), agents, Claude Code, Reinforcement Learning}
\vspace{2pt}
\skill{DevOps}{Linux, Docker, GitLab CI/CD, GitHub Actions, Jenkins, Tailscale, n8n, Kubernetes (en apprentissage)}
\vspace{2pt}
\skill{Systèmes}{C, Rust, C++, Haskell, Bash}
\vspace{2pt}
\skill{Web}{React, JavaScript, HTML/CSS}

% --- Formation
\sect{Formation}
\entry{EPITECH Lille}{Programme Grande École en informatique, 5 ans}{2024 – présent}
{\color{muted}1re et 2e années validées avec les félicitations · spécialisation visée : IA, Maths, DevOps\\
Badge Google DeepMind « AI Research Foundations » (en cours)\par}
\vspace{4pt}
\entry{Baccalauréat Général, mention Bien}{Lycée Notre-Dame des Dunes, Dunkerque}{2021 – 2024}
{\color{muted}Section européenne anglais · spécialités NSI, Mathématiques, Maths expertes\par}

% --- Langues et intérêts
\sect{Langues et centres d’intérêt}
Français (natif) · Anglais B2–C1 (TEPitech/TOEIC 910, Cambridge FIRST) · Espagnol A2\par
\vspace{2pt}
{\color{muted}Escalade, callisthénie, sport automobile, musique, dessin\par}
\vspace{2pt}
{\color{muted}Séjours : Chicago et Mexico (familles d’accueil, 2022-2023), New York (2024), Londres, Paris, Bruxelles\par}

\end{document}
